\section{Organización y personas: ir contra la naturaleza humana, seguirla y el cerebro común}

La sección anterior trató la plataforma; esta sección vuelve a la organización. Li Xiang (李想) mencionó en repetidas ocasiones la idea de «hablar siguiendo la naturaleza humana, actuar yendo contra ella». La frase puede entenderse fácilmente como una técnica de gestión, pero en el contexto de la entrevista apunta a la dificultad del cambio organizacional: las mejores prácticas, el desarrollo de capacidades, la visión de largo plazo, la revisión retrospectiva y la autocrítica suelen ir contra la naturaleza humana; en cambio, hacer lo que uno quiere, evadir en el corto plazo y justificarse a sí mismo la siguen con más facilidad. Si una organización quiere evolucionar de forma continua, debe diseñar un mecanismo que comprenda la naturaleza humana y que, al mismo tiempo, impulse un crecimiento que va en su contra.

\begin{knowledgebox}{Ir contra la naturaleza humana no es reprimir a las personas, sino combatir la inercia}
Cuando Li Xiang habla de «ir contra la naturaleza humana», se refiere a combatir la pereza, la evasión y la comodidad de corto plazo, no a negar los sentimientos de las personas. Una buena organización debe comunicarse respetando la forma de expresarse y la dignidad de las personas, pero no puede ser permisiva en el desarrollo de capacidades ni en la exigencia de resultados.
\end{knowledgebox}

\subsection{El cerebro común y el corazón común de tres a siete personas}

El apartado anterior explicó que la organización debe comprender la naturaleza humana sin consentir la inercia; este apartado examina la unidad organizacional concreta que propone Li Xiang. En la segunda mitad de la entrevista, habló de un modelo: de tres a siete personas forman un cerebro más fuerte, un corazón más fuerte y una energía más fuerte. No se trata de una discusión de grupo ordinaria, sino de reunir la complementariedad de capacidades, la energía emocional, la responsabilidad compartida y una metodología unificada. Si una persona es muy capaz pero no obtiene resultados, puede deberse a que le falta esta estructura común.

\begin{importantbox}{Cuatro condiciones del cerebro común}
Primera, los miembros se complementan entre sí en lugar de duplicarse. Segunda, comparten un objetivo y una metodología. Tercera, la relación es lo bastante cercana para dar retroalimentación sincera. Cuarta, el equipo tiene la energía para soportar la presión en conjunto, y no solo una tabla de reparto de tareas.
\end{importantbox}

\subsection{Preocuparse por los usuarios y por las personas cercanas}

Li Xiang resume la esencia de la conexión entre personas en dos formas de «preocuparse»: preocuparse por los usuarios y preocuparse por las personas cercanas. Preocuparse por los usuarios es un consenso de valores; preocuparse por las personas cercanas es la base de la colaboración. Si una organización solo atiende los asuntos y no a las personas, desgasta fácilmente las relaciones; si solo atiende las relaciones y no a los usuarios, pierde los resultados. Una buena organización debe mantener la tensión entre ambos aspectos.

\begin{warningbox}{«Juzgar los hechos y no a las personas» a veces encubre los problemas reales}
Las organizaciones de ingeniería suelen decir que se juzgan los hechos y no a las personas, lo cual evita ataques emocionales. Pero si se deja de mirar por completo a las personas, se pasan por alto su crecimiento, su motivación, su energía y la calidad de sus relaciones. Cuando Li Xiang subraya aquí «primero las personas, luego los asuntos», recuerda que los problemas de una organización se reducen, en última instancia, a si las personas se vuelven más fuertes.
\end{warningbox}

\subsection{Resumen de la sección}

La capacidad organizacional no es un diagrama de flujo, sino la manera en que las personas crecen juntas bajo una presión prolongada. En la era de la IA, las organizaciones pequeñas pueden ser más fuertes, siempre que cuenten con un cerebro común, un corazón común, una metodología común y una energía común.
